\documentclass{amsart}
\usepackage{amsmath,amssymb,amsthm,hyperref}

\newtheorem{theorem}{Theorem}
\newtheorem{lemma}{Lemma}
\newtheorem{proposition}{Proposition}
\newtheorem{corollary}{Corollary}
\theoremstyle{definition}
\newtheorem{definition}{Definition}
\newtheorem{example}{Example}
\theoremstyle{remark}
\newtheorem{remark}{Remark}

\newcommand{\zo}{[0,1]}
\newcommand{\LI}{\mathrm{LI}}
\newcommand{\EP}{\mathrm{EP}}
\newcommand{\Ax}{A}

\begin{document}

\title{The law of importation, the exchange principle, and the composition semigroup of a fuzzy implication}
\maketitle

\section{Notation and the central reformulation}

Throughout, $I\colon\zo^2\to\zo$ is a \emph{fuzzy implication}: it is non-increasing
in its first argument, non-decreasing in its second argument, and satisfies
\begin{equation}\label{eq:bd}
I(1,0)=0,\qquad I(0,x)=1,\qquad I(x,1)=1\qquad(x\in\zo).
\end{equation}
A \emph{t-norm} is a map $T\colon\zo^2\to\zo$ that is associative, commutative,
non-decreasing in each argument, and has $1$ as neutral element; recall that every
t-norm satisfies
\begin{equation}\label{eq:tnorm}
T(x,1)=T(1,x)=x,\qquad T(x,0)=0,\qquad T(x,y)\le\min(x,y)\qquad(x,y\in\zo).
\end{equation}
For a fixed $x\in\zo$ define the \emph{vertical section} (a self-map of $\zo$)
\begin{equation}\label{eq:Ax}
\Ax_x\colon\zo\to\zo,\qquad \Ax_x(z):=I(x,z).
\end{equation}
By \eqref{eq:bd} each $\Ax_x$ is non-decreasing, $\Ax_0\equiv 1$, $\Ax_x(1)=1$, and
$\Ax_1(0)=0$. The whole problem hinges on the following trivial but decisive rewriting.

\begin{lemma}[Composition reformulation]\label{lem:comp}
Let $I$ be a fuzzy implication and let $\circ$ denote composition of self-maps of $\zo$.
For every binary operation $T\colon\zo^2\to\zo$,
\[
I\ \text{satisfies}\ \LI(T)\iff \Ax_x\circ \Ax_y=\Ax_{T(x,y)}\quad\text{for all }x,y\in\zo,
\]
and
\[
I\ \text{satisfies}\ \EP\iff \Ax_x\circ \Ax_y=\Ax_y\circ \Ax_x\quad\text{for all }x,y\in\zo .
\]
\end{lemma}

\begin{proof}
For all $x,y,z$ we have, by \eqref{eq:Ax},
\[
I\bigl(x,I(y,z)\bigr)=\Ax_x\bigl(\Ax_y(z)\bigr)=(\Ax_x\circ \Ax_y)(z),\qquad
I\bigl(T(x,y),z\bigr)=\Ax_{T(x,y)}(z),
\]
and likewise $I\bigl(y,I(x,z)\bigr)=(\Ax_y\circ \Ax_x)(z)$. Two self-maps of $\zo$ are
equal iff they agree at every $z\in\zo$; reading the defining identities of
$\LI(T)$ and $\EP$ pointwise in $z$ gives both equivalences.
\end{proof}

Thus a fuzzy implication is, up to the identification $x\mapsto \Ax_x$, a family
$\mathcal{A}_I:=\{\Ax_x:x\in\zo\}$ of non-decreasing self-maps of $\zo$. The law of
importation asserts that this family is \emph{closed under composition with composition
law given by $T$}; the exchange principle asserts only that the family is
\emph{commutative} under composition. From here the comparison between the two
principles becomes transparent. We record two immediate consequences.

\begin{lemma}\label{lem:LIimpliesEP}
If $I$ satisfies $\LI(T)$ for a t-norm $T$, then $I$ satisfies $\EP$.
\end{lemma}

\begin{proof}
By Lemma~\ref{lem:comp}, $\Ax_x\circ \Ax_y=\Ax_{T(x,y)}$ and $\Ax_y\circ \Ax_x=\Ax_{T(y,x)}$
for all $x,y$. Since $T$ is commutative, $T(x,y)=T(y,x)$, hence
$\Ax_x\circ \Ax_y=\Ax_y\circ \Ax_x$, which is $\EP$.
\end{proof}

\begin{lemma}\label{lem:neutral}
If $I$ satisfies $\LI(T)$ for a t-norm $T$, then $\Ax_1$ fixes every value taken by $I$:
$I\bigl(1,I(y,z)\bigr)=I(y,z)$ for all $y,z\in\zo$. Equivalently $\Ax_1\circ \Ax_y=\Ax_y$
for all $y$.
\end{lemma}

\begin{proof}
By \eqref{eq:tnorm}, $T(1,y)=y$, so Lemma~\ref{lem:comp} gives
$\Ax_1\circ \Ax_y=\Ax_{T(1,y)}=\Ax_y$.
\end{proof}

\section{Part (a): characterization of all $I$ satisfying $\LI(T)$}

\subsection{The general characterization (arbitrary $T$)}

Lemma~\ref{lem:comp} already \emph{is} a complete characterization, but to make it a
usable statement about $I$ we phrase it intrinsically.

\begin{theorem}[Characterization of $\LI(T)$, arbitrary t-norm]\label{thm:a-general}
Let $T$ be any t-norm and $I$ a fuzzy implication. The following are equivalent.
\begin{enumerate}
\item[(i)] $I$ satisfies $\LI(T)$.
\item[(ii)] The map $\Phi_I\colon\bigl(\zo,T\bigr)\to\bigl(\mathcal{S},\circ\bigr)$,
$\Phi_I(x)=\Ax_x=I(x,\cdot)$, into the monoid $\mathcal{S}$ of self-maps of $\zo$ under
composition, is a homomorphism, i.e.
\begin{equation}\label{eq:hom}
I\bigl(T(x,y),\,\cdot\bigr)=I\bigl(x,\,I(y,\cdot)\bigr)\qquad\text{for all }x,y\in\zo .
\end{equation}
\item[(iii)] $I$ satisfies $\EP$ and, in addition, for all $x,y,z\in\zo$,
\begin{equation}\label{eq:closure}
I\bigl(x,I(y,z)\bigr)=I\bigl(T(x,y),z\bigr).
\end{equation}
\end{enumerate}
\end{theorem}

\begin{proof}
(i)$\Leftrightarrow$(ii) is Lemma~\ref{lem:comp} (condition \eqref{eq:hom} is exactly
$\Ax_{T(x,y)}=\Ax_x\circ \Ax_y$). For (i)$\Rightarrow$(iii): (i) gives \eqref{eq:closure}
by definition, and $\EP$ by Lemma~\ref{lem:LIimpliesEP}. For (iii)$\Rightarrow$(i):
\eqref{eq:closure} is the definition of $\LI(T)$.
\end{proof}

Theorem~\ref{thm:a-general} says: \emph{the fuzzy implications satisfying $\LI(T)$ are
exactly those whose vertical sections $\{I(x,\cdot)\}_{x\in\zo}$ realize $([0,1],T)$ as a
commutative semigroup of non-decreasing self-maps of $\zo$ (with $I(1,\cdot)$ playing the
role of the unit on the relevant range and $I(0,\cdot)\equiv1$ the absorbing element).}
This is the cleanest description valid for every t-norm, continuous or not.

We now make it fully explicit for continuous t-norms, using their representation theory.

\subsection{Continuous Archimedean $T$: generator form}

Recall that a continuous t-norm is Archimedean iff $T(x,x)<x$ for all $x\in(0,1)$, and
that every continuous Archimedean t-norm admits a (continuous, strictly decreasing)
\emph{additive generator} $t\colon\zo\to[0,\infty]$ with $t(1)=0$, such that
\begin{equation}\label{eq:gen}
T(x,y)=t^{-1}\!\bigl(\min\{t(0),\,t(x)+t(y)\}\bigr),
\end{equation}
where $t^{-1}$ is the pseudo-inverse; $T$ is \emph{strict} if $t(0)=\infty$ and
\emph{nilpotent} if $t(0)<\infty$. Set $e:=t(0)\in(0,\infty]$ and let
$\theta:=t^{-1}\colon[0,e]\to\zo$ be the continuous strictly decreasing bijection inverse to
$t$ (extended by $\theta(u)=0$ for $u\ge e$).

\begin{theorem}[Characterization of $\LI(T)$, continuous Archimedean $T$]\label{thm:a-arch}
Let $T$ be a continuous Archimedean t-norm with additive generator $t$, $e=t(0)$, and let
$I$ be a fuzzy implication. Define
\begin{equation}\label{eq:J}
J\colon[0,e]\times\zo\to\zo,\qquad J(u,z):=I\bigl(\theta(u),z\bigr)=I\bigl(t^{-1}(u),z\bigr).
\end{equation}
Then $I$ satisfies $\LI(T)$ if and only if $J$ is a \emph{truncated translation
semigroup action}: for all $u,v\in[0,e]$ and all $z\in\zo$,
\begin{equation}\label{eq:flow}
J\bigl(u,J(v,z)\bigr)=J\bigl(\min\{u+v,\;e\},\,z\bigr).
\end{equation}
\end{theorem}

\begin{proof}
Because $t$ is a continuous decreasing bijection of $\zo$ onto $[0,e]$, the substitution
$x=\theta(u)$, $y=\theta(v)$ is a bijection of $\zo^2$ onto $[0,e]^2$. Under it,
$\LI(T)$, i.e. $I(x,I(y,z))=I(T(x,y),z)$, becomes
\[
J\bigl(u,J(v,z)\bigr)=I\bigl(\theta(u),I(\theta(v),z)\bigr)
=I\bigl(T(\theta(u),\theta(v)),z\bigr).
\]
By \eqref{eq:gen} and $t(\theta(u))=u$ (for $u\le e$), we have
$T(\theta(u),\theta(v))=t^{-1}(\min\{e,u+v\})=\theta(\min\{e,u+v\})$, so the right-hand
side equals $J(\min\{u+v,e\},z)$. Hence $\LI(T)$ is equivalent to \eqref{eq:flow} for all
$u,v\in[0,e]$, $z\in\zo$.
\end{proof}

Equation~\eqref{eq:flow} states precisely that the family of non-decreasing self-maps
$\{J(u,\cdot):u\in[0,e]\}$ is a one-parameter semigroup under composition, additively
parametrized and truncated at the absorbing level $e$ (where $J(e,z)=I(0,z)=1$). For the
\emph{strict} case ($e=\infty$) this is a genuine one-parameter flow
$\{J(u,\cdot)\}_{u\ge0}$ with $J(0,\cdot)=\mathrm{id}$ on the relevant range and
$J(u,\cdot)\to 1$ pointwise as $u\to\infty$.

\begin{example}[Strict case]\label{ex:strict}
Take $T=T_{\mathbf P}$ (product), with generator $t(x)=-\ln x$, $e=\infty$,
$\theta(u)=e^{-u}$. The Yager implication $I(x,y)=y^{x}$ (with $0^0:=1$) gives
$J(u,z)=z^{e^{-u}}$, and indeed
$J(u,J(v,z))=\bigl(z^{e^{-v}}\bigr)^{e^{-u}}=z^{e^{-(u+v)}}=J(u+v,z)$, confirming
\eqref{eq:flow}. Equivalently $I(x,I(y,z))=z^{xy}=I(xy,z)=I(T_{\mathbf P}(x,y),z)$.
\end{example}

\begin{example}[Nilpotent case]\label{ex:nilp}
Take $T=T_{\mathbf L}$ (\L ukasiewicz), $T_{\mathbf L}(x,y)=\max\{0,x+y-1\}$, with
generator $t(x)=1-x$, $e=1$, $\theta(u)=1-u$. The \L ukasiewicz implication
$I(x,y)=\min\{1,1-x+y\}$ gives $J(u,z)=\min\{1,u+z\}$, and
$J(u,J(v,z))=\min\{1,u+\min\{1,v+z\}\}=\min\{1,u+v+z\}=J(\min\{u+v,1\},z)$,
confirming \eqref{eq:flow}.
\end{example}

\subsection{General continuous $T$: ordinal-sum form}

A general continuous t-norm is an \emph{ordinal sum}
$T=\bigl(\langle a_k,b_k,T_k\rangle\bigr)_{k\in K}$ of continuous Archimedean summands
$T_k$ on a countable family of pairwise disjoint open subintervals
$(a_k,b_k)\subseteq\zo$; outside the squares $[a_k,b_k]^2$ one has $T(x,y)=\min(x,y)$.
Theorem~\ref{thm:a-general} reduces $\LI(T)$ to the homomorphism property
$\Ax_{T(x,y)}=\Ax_x\circ\Ax_y$, which splits according to the ordinal-sum structure:
\begin{itemize}
\item Whenever $(x,y)$ lies outside all summand squares, $T(x,y)=\min(x,y)$, so $\LI(T)$
requires there $I(\min(x,y),\cdot)=I(x,\cdot)\circ I(y,\cdot)$; in particular, on the
idempotent set $E:=\{x:T(x,x)=x\}=\zo\setminus\bigcup_k(a_k,b_k)$ each $\Ax_x$ is
idempotent, $\Ax_x\circ\Ax_x=\Ax_x$, and the family $\{\Ax_x:x\in E\}$ is a chain of
commuting idempotents composing by $\min$ (the "G\"odel pattern").
\item On each summand square $[a_k,b_k]^2$, $T$ acts as the Archimedean $T_k$, so by the
argument of Theorem~\ref{thm:a-arch} the restricted family
$\{I(\theta_k(u),\cdot):u\in[0,e_k]\}$ (with $\theta_k$ the generator-inverse of $T_k$)
must be a truncated translation semigroup as in \eqref{eq:flow}, the truncation level
$e_k=t_k(a_k)$ corresponding to collapse onto $\Ax_{a_k}$.
\end{itemize}
Conversely, if these two patterns hold and are compatible at the endpoints $a_k,b_k$
(continuity of $T$ forces $T(a_k,y)=\min(a_k,y)=a_k$ for $y\ge a_k$, hence
$\Ax_{a_k}\circ\Ax_y=\Ax_{a_k}$ there), then $\Ax_{T(x,y)}=\Ax_x\circ\Ax_y$ holds for all
$x,y$ and $\LI(T)$ follows from Theorem~\ref{thm:a-general}. This is the complete
description for an arbitrary continuous t-norm: \emph{$\LI(T)$ holds iff the vertical
sections of $I$ form, summand by summand, an additively parametrized semigroup of maps,
glued along the idempotent skeleton by $\min$.}

\begin{example}[$T=\min$]\label{ex:min}
Here $E=\zo$ (every point idempotent), so $\LI(\min)$ holds iff
$I(\min(x,y),\cdot)=I(x,\cdot)\circ I(y,\cdot)$ for all $x,y$; this is exactly $\EP$
\emph{together with} the requirement that the sections form a $\min$-chain of commuting
maps. The G\"odel implication $I(x,y)=1$ if $x\le y$, $=y$ otherwise, satisfies
$I(x,I(y,z))=I(\min(x,y),z)$, hence $\LI(\min)$.
\end{example}

\section{Part (b): the converse, and uniqueness}

\subsection{The converse fails in general}

\begin{theorem}[Counterexample]\label{thm:counter}
There is a fuzzy implication $I^{\ast}$ that satisfies the exchange principle but satisfies
$\LI(T)$ for no t-norm $T$ whatsoever.
\end{theorem}

\begin{proof}
Define $w\colon\zo\to\{0,1,2\}$ by
\[
w(x)=\begin{cases}0,& x=0,\\ 1,& 0<x\le\tfrac12,\\ 2,& \tfrac12<x\le1,\end{cases}
\]
and set, with the convention $0^{0}:=1$,
\begin{equation}\label{eq:Istar}
I^{\ast}(x,y):=y^{\,w(x)}\qquad(x,y\in\zo).
\end{equation}

\emph{$I^{\ast}$ is a fuzzy implication.} For fixed $y\in(0,1)$, $y^{w}$ is decreasing in
$w$, and $w$ is non-decreasing, so $x\mapsto I^{\ast}(x,y)=y^{w(x)}$ is non-increasing;
for $y\in\{0,1\}$ the section is constant ($0^{w}$: equals $1$ if $w=0$ and $0$ if $w>0$,
still non-increasing in $x$; $1^{w}\equiv1$). For fixed $x$, $y\mapsto y^{w(x)}$ is
non-decreasing on $\zo$ (as $w(x)\ge0$). Finally $I^{\ast}(1,0)=0^{2}=0$,
$I^{\ast}(0,y)=y^{0}=1$, and $I^{\ast}(x,1)=1^{w(x)}=1$, so \eqref{eq:bd} holds.

\emph{$I^{\ast}$ satisfies $\EP$.} For all $x,y,z$,
\[
I^{\ast}\bigl(x,I^{\ast}(y,z)\bigr)=\bigl(z^{w(y)}\bigr)^{w(x)}=z^{\,w(x)w(y)},
\]
which is symmetric in $x,y$ because $w(x)w(y)=w(y)w(x)$ in $\mathbb{R}$; hence
$I^{\ast}(x,I^{\ast}(y,z))=I^{\ast}(y,I^{\ast}(x,z))$.

\emph{$I^{\ast}$ satisfies $\LI(T)$ for no t-norm $T$.} Suppose, for contradiction, that
$\LI(T)$ held for some t-norm $T$. Fix any $z\in(0,1)$. Then for all $x,y$,
\[
z^{\,w(x)w(y)}=I^{\ast}\bigl(x,I^{\ast}(y,z)\bigr)=I^{\ast}\bigl(T(x,y),z\bigr)
=z^{\,w(T(x,y))}.
\]
Since $z\in(0,1)$, the map $a\mapsto z^{a}$ is injective on $[0,\infty)$, so
\begin{equation}\label{eq:wmult}
w(x)\,w(y)=w\bigl(T(x,y)\bigr)\qquad\text{for all }x,y\in\zo.
\end{equation}
Take $x=y=1$. Because $1$ is the neutral element of $T$, $T(1,1)=1$, so the right side of
\eqref{eq:wmult} is $w(1)=2$, while the left side is $w(1)^{2}=4$. As $4\ne2$, this is a
contradiction. Hence no t-norm $T$ exists with $\LI(T)$.
\end{proof}

\begin{remark}
The obstruction is structural and is read directly from Lemma~\ref{lem:comp}: the vertical
sections of $I^{\ast}$ are $\Ax_x=(\cdot)^{w(x)}$, and $\Ax_x\circ\Ax_y=(\cdot)^{w(x)w(y)}$.
The family $\{\Ax_x\}=\{\mathrm{id}^{0},\mathrm{id}^{1},\mathrm{id}^{2}\}$ is commutative
(so $\EP$ holds) but is \emph{not closed under composition}: composing the exponent-$2$ map
with itself produces the exponent-$4$ map, which is not in the family. By
Lemma~\ref{lem:comp} no closure law $\Ax_x\circ\Ax_y=\Ax_{T(x,y)}$ can exist, let alone
one given by a t-norm.
\end{remark}

\subsection{Exactly when the converse holds}

The counterexample isolates the missing ingredient: beyond commutativity ($\EP$) one needs
\emph{closure under composition}. We now turn this into an exact characterization.
Introduce the section relation
\[
x\sim x'\ :\Longleftrightarrow\ I(x,\cdot)=I(x',\cdot),\qquad\text{i.e. }\Ax_x=\Ax_{x'} .
\]

\begin{theorem}[Characterization of the converse]\label{thm:b-char}
Let $I$ be a fuzzy implication satisfying $\EP$. Then $I$ satisfies $\LI(T)$ for some
t-norm $T$ if and only if all four of the following hold.
\begin{enumerate}
\item[(C1)] \emph{Composition closure:} for all $x,y\in\zo$ there exists $s\in\zo$ with
$\Ax_x\circ\Ax_y=\Ax_s$, i.e. $I(x,I(y,\cdot))=I(s,\cdot)$.
\item[(C2)] \emph{Left neutrality on the section algebra:} $\Ax_1=I(1,\cdot)$ acts as the
identity on the range of every $\Ax_y$, equivalently $\Ax_1\circ\Ax_y=\Ax_y$ for all $y$.
\item[(C3)] \emph{Monotone selectability:} the index $s$ in $(C1)$ can be chosen so that the
resulting operation $T$ defined by $\Ax_{T(x,y)}=\Ax_x\circ\Ax_y$ is non-decreasing in each
argument and satisfies $T(x,0)=0$, $T(x,1)=x$.
\item[(C4)] \emph{(Always available; included for completeness)} $\EP$ holds, giving the
commutativity of $T$.
\end{enumerate}
Moreover, when $I$ satisfies $\LI(T)$, the t-norm $T$ is \textbf{uniquely determined by $I$}
if and only if $\sim$ is trivial on the relevant index set, i.e. iff the map
$x\mapsto I(x,\cdot)$ is injective; in that case $T$ is forced by
\begin{equation}\label{eq:Tforced}
\Ax_{T(x,y)}=\Ax_x\circ\Ax_y .
\end{equation}
\end{theorem}

\begin{proof}
\emph{Necessity.} Assume $\LI(T)$ for a t-norm $T$. By Lemma~\ref{lem:comp},
$\Ax_x\circ\Ax_y=\Ax_{T(x,y)}$, giving (C1) with $s=T(x,y)$. (C2) is
Lemma~\ref{lem:neutral}. For (C3): with this same choice $s=T(x,y)$, monotonicity of $T$
and the boundary values $T(x,0)=0$, $T(x,1)=x$ hold by \eqref{eq:tnorm}, and indeed
$\Ax_{T(x,y)}=\Ax_x\circ\Ax_y$ holds by construction. (C4) is Lemma~\ref{lem:LIimpliesEP}.

\emph{Sufficiency.} Assume (C1)--(C4). Let $T\colon\zo^2\to\zo$ be the operation furnished by
(C3): for each $(x,y)$ choose $T(x,y)$ with $\Ax_{T(x,y)}=\Ax_x\circ\Ax_y$ (possible by (C1)),
the choice being made so that (C3) holds. We verify $T$ is a t-norm and that $\LI(T)$ holds.

\textit{$\LI(T)$.} By the defining property of $T$, $\Ax_x\circ\Ax_y=\Ax_{T(x,y)}$ for all
$x,y$; Lemma~\ref{lem:comp} gives $\LI(T)$.

\textit{Neutral element $1$.} By (C2), $\Ax_1\circ\Ax_y=\Ax_y=\Ax_{T(1,y)}$, and by (C3)
$T(1,y)=y$ (we may, and do, take $T(1,y):=y$ since $\Ax_{y}=\Ax_1\circ\Ax_y$); symmetrically
$T(x,1)=x$. Hence $1$ is neutral.

\textit{Commutativity.} By $\EP$ (C4) and Lemma~\ref{lem:comp},
$\Ax_x\circ\Ax_y=\Ax_y\circ\Ax_x$, so $\Ax_{T(x,y)}=\Ax_{T(y,x)}$; choosing the same
representative for equal section-maps (which is consistent since both equal the same
self-map) we get $T(x,y)=T(y,x)$.

\textit{Associativity.} For all $x,y,u$,
\[
\Ax_{T(T(x,y),u)}=\Ax_{T(x,y)}\circ\Ax_u=(\Ax_x\circ\Ax_y)\circ\Ax_u
=\Ax_x\circ(\Ax_y\circ\Ax_u)=\Ax_x\circ\Ax_{T(y,u)}=\Ax_{T(x,T(y,u))},
\]
using associativity of composition of self-maps. Thus $T(T(x,y),u)$ and $T(x,T(y,u))$ index
the same section-map; choosing representatives consistently gives
$T(T(x,y),u)=T(x,T(y,u))$.

\textit{Monotonicity.} This is (C3).

Hence $T$ is a commutative, associative, monotone binary operation with neutral element $1$,
i.e. a t-norm, and $I$ satisfies $\LI(T)$.

\emph{Uniqueness.} Suppose $I$ satisfies both $\LI(T)$ and $\LI(T')$. By
Lemma~\ref{lem:comp},
\[
\Ax_{T(x,y)}=\Ax_x\circ\Ax_y=\Ax_{T'(x,y)}\qquad(x,y\in\zo),
\]
so $T(x,y)\sim T'(x,y)$ for all $x,y$. If $x\mapsto\Ax_x$ is injective, $\sim$ is trivial and
$T(x,y)=T'(x,y)$, proving uniqueness and \eqref{eq:Tforced}.

Conversely, if $x\mapsto\Ax_x$ is not injective then $T$ need not be unique; this is not an
abstract possibility but is realized concretely in Theorem~\ref{thm:nonunique}, where a single
$\EP$ implication with non-injective section map admits the two distinct t-norms $\min$ and the
product. Hence injectivity of $x\mapsto\Ax_x$ is precisely the condition separating the unique
and the non-unique regimes: it is sufficient for uniqueness (just shown) and, by
Theorem~\ref{thm:nonunique}, its failure can destroy uniqueness.
\end{proof}

\begin{remark}[Reading of Theorem~\ref{thm:b-char}]
The decisive extra hypothesis beyond $\EP$ is \emph{(C1), closure of the vertical sections
under composition}; (C2) is the (mild) demand that $I(1,\cdot)$ be neutral, automatically
true e.g. when $I$ satisfies the left-neutrality property $I(1,y)=y$; (C3) is a compatibility
of the selection with the order, automatic whenever the sections are linearly ordered by the
pointwise order and indexed monotonically. In the well-behaved sub-class --- fuzzy
implications that satisfy $\EP$, the left-neutrality $I(1,y)=y$, and for which
$x\mapsto I(x,\cdot)$ is injective --- conditions (C1)--(C4) reduce to the single requirement
(C1), and then $T$ exists and is \textbf{unique}, given by \eqref{eq:Tforced}. This is the
precise class for which "$\EP\Rightarrow\LI(T)$ for a unique $T$" holds.
\end{remark}

\subsection{Uniqueness fails without injectivity}

We make the uniqueness criterion concrete.

\begin{theorem}[Non-uniqueness]\label{thm:nonunique}
There is a fuzzy implication $I^{\circ}$ satisfying $\EP$ that satisfies $\LI(T)$ for
\emph{every} t-norm $T$ without zero divisors (a \emph{positive} t-norm, i.e.
$T(x,y)=0\Rightarrow x=0$ or $y=0$), and for no t-norm with zero divisors. In particular
$I^{\circ}$ satisfies $\LI(T_{\mathbf M})$ (minimum) and $\LI(T_{\mathbf P})$ (product),
two distinct t-norms; hence $T$ is not uniquely determined by $I^{\circ}$.
\end{theorem}

\begin{proof}
Define
\[
I^{\circ}(x,y):=\begin{cases}1,& x=0\ \text{or}\ y=1,\\[1mm] 0,& x>0\ \text{and}\ y=0,\\[1mm]
y,& x>0\ \text{and}\ 0<y<1.\end{cases}
\]

\emph{$I^{\circ}$ is a fuzzy implication.} For fixed $y\in(0,1)$ the section in $x$ is
$1$ at $x=0$ and $y$ for $x>0$: non-increasing. For $y=0$ it is $1$ at $x=0$ and $0$ for
$x>0$: non-increasing. For $y=1$ it is constantly $1$. For fixed $x>0$, the section in $y$ is
$0,\,y,\,1$ on $\{0\},(0,1),\{1\}$ respectively: non-decreasing; for $x=0$ it is $1$. Boundary
conditions: $I^{\circ}(1,0)=0$, $I^{\circ}(0,y)=1$, $I^{\circ}(x,1)=1$, so \eqref{eq:bd} holds.

\emph{Vertical sections.} $\Ax_0\equiv1$, and for every $x>0$,
\[
\Ax_x(z)=I^{\circ}(x,z)=\begin{cases}0,& z=0,\\ z,& 0<z<1,\\ 1,& z=1,\end{cases}
\]
which is one and the same self-map $g$, \emph{independent of} $x>0$. Thus
$\Ax_x=g$ for all $x\in(0,1]$ and $\Ax_0\equiv1$. The map $x\mapsto\Ax_x$ takes only two
values and is far from injective: every $x>0$ has $\Ax_x=g$.

\emph{Composition.} Note $g\circ g=g$ (since $g$ is the identity on $(0,1)$ and fixes $0,1$),
$g\circ\mathbf 1=\mathbf 1$ and $\mathbf 1\circ g=\mathbf 1$, $\mathbf 1\circ\mathbf 1=\mathbf
1$, where $\mathbf 1\equiv1=\Ax_0$. Hence for a candidate $T$ the law
$\Ax_x\circ\Ax_y=\Ax_{T(x,y)}$ reads, by Lemma~\ref{lem:comp}:
\[
\begin{aligned}
&x=0\ \text{or}\ y=0:\quad \Ax_{T(x,y)}=\mathbf 1\iff T(x,y)=0\ \text{(use $\Ax_0=\mathbf1$,
$\Ax_s=g\neq\mathbf1$ for $s>0$)};\\
&x>0\ \text{and}\ y>0:\quad \Ax_{T(x,y)}=g\circ g=g\iff T(x,y)>0 .
\end{aligned}
\]
For any t-norm, $x=0$ or $y=0$ forces $T(x,y)=0$ by \eqref{eq:tnorm}, so the first line
always holds. Therefore
\[
I^{\circ}\ \text{satisfies}\ \LI(T)\iff
\bigl[\,x>0\ \text{and}\ y>0\ \Rightarrow\ T(x,y)>0\,\bigr]\iff T\ \text{has no zero divisors}.
\]
Both $T_{\mathbf M}=\min$ and $T_{\mathbf P}(x,y)=xy$ are positive (no zero divisors), so
$I^{\circ}$ satisfies $\LI(\min)$ and $\LI(T_{\mathbf P})$; the \L ukasiewicz t-norm has zero
divisors ($T_{\mathbf L}(\tfrac12,\tfrac12)=0$) and is excluded. Finally
$I^{\circ}$ satisfies $\EP$ by Theorem~\ref{thm:b-char} (it satisfies $\LI(\min)$, hence
$\EP$ by Lemma~\ref{lem:LIimpliesEP}), or directly since
$I^{\circ}(x,I^{\circ}(y,z))=I^{\circ}(y,I^{\circ}(x,z))$ both equal $g(g(z))=g(z)$ when
$x,y>0$ and equal $1$ as soon as $x=0$ or $y=0$.

Since $\min\neq T_{\mathbf P}$ are both admissible, $T$ is not uniquely determined.
\end{proof}

\begin{remark}
Theorem~\ref{thm:nonunique} exhibits exactly the failure mode predicted by
Theorem~\ref{thm:b-char}: $x\mapsto I^{\circ}(x,\cdot)$ is non-injective (constant on
$(0,1]$), so the index $T(x,y)$ of $\Ax_x\circ\Ax_y$ is only determined up to $\sim$ --- here
only "$T(x,y)>0$ or $=0$" is forced --- and any positive t-norm fills in the rest.
\end{remark}

\section{Summary of answers}

\textbf{(a)} For an arbitrary t-norm $T$, a fuzzy implication $I$ satisfies $\LI(T)$ if and
only if its vertical sections $\Ax_x=I(x,\cdot)$ form a commutative semigroup of
non-decreasing self-maps of $\zo$ whose composition law is $T$:
\[
I(T(x,y),\cdot)=I(x,I(y,\cdot))\qquad\text{for all }x,y
\]
(Theorem~\ref{thm:a-general}). For a continuous Archimedean $T$ with additive generator $t$
and $e=t(0)$, this is equivalent to the truncated translation-semigroup law
$J(u,J(v,z))=J(\min\{u+v,e\},z)$ for $J(u,z)=I(t^{-1}(u),z)$ (Theorem~\ref{thm:a-arch}); for a
general continuous $T$, it is the ordinal-sum gluing of such semigroups along the idempotent
skeleton, with the $\min$-pattern off the summand squares (\S2.3).

\textbf{(b)} The converse is \emph{false}: there is an $\EP$ implication satisfying $\LI(T)$
for \emph{no} t-norm (Theorem~\ref{thm:counter}, $I^{\ast}(x,y)=y^{w(x)}$ with the three-valued
exponent $w$). Precisely, an $\EP$ implication satisfies $\LI(T)$ for some t-norm iff its
vertical sections are closed under composition (plus the neutrality and monotone-selection
conditions (C2),(C3) of Theorem~\ref{thm:b-char}); the commutativity needed is supplied by
$\EP$. When such a $T$ exists, injectivity of $x\mapsto I(x,\cdot)$ \emph{forces} uniqueness, with
$T$ given by $\Ax_{T(x,y)}=\Ax_x\circ\Ax_y$; if this map is not injective,
$T$ need not be unique --- e.g. $I^{\circ}$ of Theorem~\ref{thm:nonunique} satisfies $\LI(T)$
for every positive t-norm, in particular for both $\min$ and the product.

\end{document}
